\Needspace{15\baselineskip}\Article{9}

Secondly, we shall consider the series

\[
\begin{aligned}
&1+x^{\frac12}\frac{1-x^m}{1-x}+x\frac{(1-x^m)(1-x^{m-1})}{(1-x)(1-xx)}\\
&\qquad+x^{\frac32}\frac{(1-x^m)(1-x^{m-1})(1-x^{m-2})}{(1-x)(1-xx)(1-x^3)}+\text{etc.}
\end{aligned}
\]

or

\[
\begin{aligned}
1+x^{\frac12}(m,1)+x(m,2)+x^{\frac32}(m,3)+xx(m,4)+\text{etc.}
\end{aligned}
\]

which we shall denote by $F(x,m)$.\quad We restrict this investigation to the case where $m$ is a positive integer, so that this series too always ends
\SourceLeaf{31}
at its term in position $m+1$ (with $=x^{\frac12m}(m,m)$).\quad Since

\[
\begin{aligned}
(m,m)=1,\quad(m,m-1)=(m,1),\quad(m,m-2)=(m,2),\;\text{etc.}
\end{aligned}
\]

the series can also be written as follows:

\[
\begin{aligned}
F(x,m)&=x^{\frac12m}+x^{\frac12(m-1)}(m,1)\\&\qquad+x^{\frac12(m-2)}(m,2)+x^{\frac12(m-3)}(m,3)+\text{etc.}
\end{aligned}
\]

This gives

\[
\begin{aligned}
(1+x^{\frac12m+\frac12})F(x,m)&=1+x^{\frac12}(m,1)+x(m,2)\\&\qquad+x^{\frac32}(m,3)+\text{etc.}\\
&+x^{\frac12}\pdot x^m+x\pdot x^{m-1}(m,1)\\&\qquad+x^{\frac32}\pdot x^{m-2}(m,2)+\text{etc.}
\end{aligned}
\]

Since we have (Article 5, II)

\[
\begin{aligned}
(m,1)+x^m&=(m+1,1)\\
(m,2)+x^{m-1}(m,1)&=(m+1,2)\\
(m,3)+x^{m-2}(m,2)&=(m+1,3),\;\text{etc.}
\end{aligned}
\]

we obtain

\[
\begin{aligned}
(1+x^{\frac12m+\frac12})F(x,m)=F(x,m+1)\ .\; .\; .\; .\; .\; .\; .\ [3]
\end{aligned}
\]

But $F(x,0)=1$, so that

\[
\begin{aligned}
F(x,1)&=1+x^{\frac12}\\
F(x,2)&=(1+x^{\frac12})(1+x)\\
F(x,3)&=(1+x^{\frac12})(1+x)(1+x^{\frac32}),\;\text{etc.}
\end{aligned}
\]

or, in general,

\[
\begin{aligned}
F(x,m)&=(1+x^{\frac12})(1+x)(1+x^{\frac32})\\&\qquad\dotsfour(1+x^{\frac12m})\ .\; .\; .\; .\quad[4]
\end{aligned}
\]

\Needspace{12\baselineskip}\Article{10}

With these preliminary investigations in place, let us now come closer to our aim.\quad For a prime value of $n$, the squares $1$, $4$, $9\dotsfour(\frac12(n-1))^2$ are all incongruent to one another modulo $n$; their least residues for this modulus must therefore be identical to the numbers $a$, and hence

\[
\begin{aligned}
\Sigma\cos ak\omega&=\cos k\omega+\cos4k\omega+\cos9k\omega\\&\qquad+\text{etc.}+\cos(\tfrac12(n-1))^2k\omega\\
\Sigma\sin ak\omega&=\sin k\omega+\sin4k\omega+\sin9k\omega\\&\qquad+\text{etc.}+\sin(\tfrac12(n-1))^2k\omega
\end{aligned}
\]

Likewise, since the same squares $1$, $4$, $9\dotsfour(\frac12(n-1))^2$, in reverse order, are congruent to $(\frac12(n+1))^2$, $(\frac12(n+3))^2$, $(\frac12(n+5))^2\dotsfive(n-1)^2$, we also have

\SourceLeaf{32}

\[
\begin{aligned}
\Sigma\cos ak\omega&=\cos(\tfrac12(n+1))^2k\omega\\&\qquad+\cos(\tfrac12(n+3))^2k\omega+\text{etc.}+\cos(n-1)^2k\omega\\
\Sigma\sin ak\omega&=\sin(\tfrac12(n+1))^2k\omega\\&\qquad+\sin(\tfrac12(n+3))^2k\omega+\text{etc.}+\sin(n-1)^2k\omega
\end{aligned}
\]

\Needspace{9\baselineskip}Thus, putting

\[
\begin{aligned}
T&=1+\cos k\omega+\cos4k\omega+\cos9k\omega\\&\qquad+\text{etc.}+\cos(n-1)^2k\omega\\
U&=\phantom{1+}\sin k\omega+\sin4k\omega+\sin9k\omega\\&\qquad+\text{etc.}+\sin(n-1)^2k\omega
\end{aligned}
\]

we obtain

\[
\begin{aligned}
1+2\Sigma\cos ak\omega&=T\\
2\Sigma\sin ak\omega&=U
\end{aligned}
\]

This shows that the summations proposed in Article 1 depend on summing the series $T$ and $U$. We shall therefore leave the former aside and direct our investigation to these, carrying it out generally enough to include both prime values of $n$ and any composite values.\quad We shall, however, suppose that the number $k$ is relatively prime to $n$: the case where $k$ and $n$ have a common divisor can be reduced to this one without difficulty.

\Needspace{12\baselineskip}\Article{11}

Let us denote the imaginary quantity $\surd{-1}$ by $i$ and put

\[
\begin{aligned}
\cos k\omega+i\sin k\omega=r
\end{aligned}
\]

Then $r^n=1$, or $r$ is a root of the equation $r^n-1=0$.\quad It is easy to see that all the numbers $k$, $2k$, $3k\dotsfour(n-1)k$ are not divisible by $n$ and are incongruent to one another modulo $n$. Thus the powers of $r$

\[
\begin{aligned}
1,\;r,\;rr,\;r^3\dotsfour r^{n-1}
\end{aligned}
\]

are all unequal, while each also satisfies the equation $x^n-1=0$. For this reason these powers represent \emph{all} the roots of the equation $x^n-1=0$.

These conclusions would not hold if $k$ had a common divisor with $n$. For if $\nu$ were such a common divisor, $k\pdot\dfrac n\nu$ would be divisible by $n$, and hence a power lower than $r^n$, namely $r^{\frac n\nu}$, would equal unity.\quad In this case, therefore, the powers of $r$ represent at most $\dfrac n\nu$ roots of the equation $x^n-1=0$, and indeed give that many distinct roots if $\nu$ is the \emph{greatest} common divisor of the numbers\SourceLeaf{33}
$k$, $n$.\quad In our case, where $k$ and $n$ are supposed relatively prime, $r$ can conveniently be called a \emph{proper root} of the equation $x^n-1=0$. In the other case, where $k$ and $n$ have a (greatest) common divisor $\nu$, $r$ would be called an \emph{improper root} of that equation. Then, however, the same $r$ would evidently be a proper root of the equation $x^{\frac n\nu}-1=0$.\quad The simplest improper root is unity, and when $n$ is a prime number there are no other improper roots at all.

\Needspace{13\baselineskip}\Article{12}

Now, if we put

\[
\begin{aligned}
W=1+r+r^4+r^9+\text{etc.}+\text{etc.}+r^{(n-1)^2}
\end{aligned}
\]

we evidently have $W=T+iU$, so that $T$ is the real part of $W$, and $U$ is obtained from the imaginary part of $W$ by removing the factor $i$.\quad The whole question is therefore reduced to finding the sum $W$. For this purpose we can use either the series considered in Article 6 or the one we showed how to sum in Article 9; the former, however, is less suitable when $n$ is even.\quad Nevertheless, we hope that readers will welcome our treating the case where $n$ is odd by two methods.

First, then, suppose that $n$ is an odd number, that $r$ denotes any proper root of the equation $x^n-1=0$, and that in the function $f(x,m)$ we put $x=r$ and $m=n-1$.\quad We evidently obtain

\[
\begin{aligned}
\frac{1-x^m}{1-x}&=\frac{1-r^{-1}}{1-r}=-r^{-1}\\
\frac{1-x^{m-1}}{1-xx}&=\frac{1-r^{-2}}{1-rr}=-r^{-2}\\
\frac{1-x^{m-2}}{1-x^3}&=\frac{1-r^{-3}}{1-r^3}=-r^{-3},\;\text{etc.}
\end{aligned}
\]

as far as

\[
\begin{aligned}
\frac{1-x}{1-x^m}=\frac{1-r^{-m}}{1-r^m}=-r^{-m}
\end{aligned}
\]

(It is worth pointing out that these equations hold only insofar as $r$ is supposed to be a proper root: if $r$ were an improper root, the numerator and denominator of some of these fractions would vanish together, and the fractions would therefore become indeterminate).

From this we deduce the following equation

\SourceLeaf{34}

\[
\begin{aligned}
f(r,n-1)&=1+r^{-1}+r^{-3}+r^{-6}+\text{etc.}+r^{-\frac12(n-1)n}\\
&=(1-r)(1-r^3)(1-r^5)\dotsfive(1-r^{n-2})
\end{aligned}
\]

The same equation still holds if we replace $r$ by $r^\lambda$, where $\lambda$ denotes any integer relatively prime to $n$: for then $r^\lambda$ too is a proper root of the equation $x^n-1=0$.\quad Let us therefore replace $r$ by $r^{n-2}$, or equivalently $r^{-2}$. We obtain

\[
\begin{aligned}
&1+r^2+r^6+r^{12}+\text{etc.}+r^{(n-1)n}\\&\qquad=(1-r^{-2})(1-r^{-6})(1-r^{-10})\ldots(1-r^{-2(n-2)})
\end{aligned}
\]

Let us multiply both sides of this equation by

\[
\begin{aligned}
r\pdot r^3\pdot r^5\dotsfour r^{(n-2)}=r^{\frac14(n-1)^2}
\end{aligned}
\]

\Needspace{13\baselineskip}and, since

\[
\begin{aligned}
r^{2+\frac14(n-1)^2}&=r^{\frac14(n-3)^2},\\r^{(n-1)n+\frac14\,n-1)^2}=r^{\frac14(n+1)^2}\\
r^{6+\frac14(n-1)^2}&=r^{\frac14(n-5)^2},\\r^{(n-2)(n-1)+\frac14(n-1)^2}=r^{\frac14(n+3)^2}\\
r^{12+\frac14(n-1)^2}&=r^{\frac14(n-7)^2},\\r^{(n-3)(n-2)+\frac14(n-1)^2}=r^{\frac14(n+5)^2},\;\text{etc.}
\end{aligned}
\]

the following equation results:

\[
\begin{aligned}
&r^{\frac14(n-1)^2}+r^{\frac14(n-3)^2}+r^{\frac14(n-5)^2}+\text{etc.}+r+1\\
&+r^{\frac14(n+1)^2}+r^{\frac14(n+3)^2}+r^{\frac14(n+5)^2}\\&\qquad+\text{etc.}+r^{\frac14(2n-2)^2}\\
&\qquad=(r-r^{-1})(r^3-r^{-3})(r^5-r^{-5})\dotsfour(r^{n-2}-r^{-n+2})
\end{aligned}
\]

\Needspace{7\baselineskip}or, arranging the terms on the left side differently,

\[
\begin{aligned}
&1+r+r^4+\text{etc.}+r^{(n-1)^2}\\&\qquad=(r-r^{-1})(r^3-r^{-3})\ldots(r^{n-2}-r^{-n+2})\quad\dotstwo\quad[5]
\end{aligned}
\]
